\documentclass[aps,prd,reprint,amsmath,amssymb,nofootinbib,longbibliography,superscriptaddress]{revtex4-2}
\usepackage{booktabs}
\usepackage{xcolor}
\usepackage{hyperref}
\hypersetup{colorlinks=true,linkcolor=blue!50!black,citecolor=blue!50!black,urlcolor=blue!50!black}
\usepackage{bm}

\newcommand{\Msun}{M_\odot}
\newcommand{\kHz}{\,\mathrm{kHz}}
\newcommand{\cs}{c_s}
\newcommand{\BDNK}{BDNK}
\newcommand{\void}{\textsc{void}}
\newcommand{\de}{\delta\varepsilon}
\newcommand{\dS}{\delta S}
\newcommand{\dv}{\delta v}
\newcommand{\dpp}{\delta p}
\newcommand{\sg}{\sqrt{\gamma}}
\newcommand{\KO}{\sigma_{\rm KO}}

\begin{document}

\title{A cut-cell Cartesian method for linear and causal-viscous oscillations of relativistic
stars in the Cowling approximation}

\author{N. Fernandes}
\affiliation{California Institute of Technology, Pasadena, CA 91125, USA}

\date{\today}

\begin{abstract}
We present a 3+1D time-domain method for the linear oscillations of a relativistic star on a
frozen Tolman--Oppenheimer--Volkoff background (the Cowling approximation), on a Cartesian grid
that does not conform to the stellar surface. The linearized fluid equations are evolved in a
covariant-momentum form in which the factor $(\varepsilon_0+p_0)$ cancels from every right-hand
side, removing the surface singularity of the velocity formulation; the lapse multiplying the
Euler equation is derived from $\nabla_\mu\delta T^{\mu}{}_\nu = 0$ and shown to be essential,
its omission producing a resolution-independent $+5\%$ frequency error. The stellar surface is
treated by cut cells: the volume fraction and the six face apertures of every surface cell reduce
to one-dimensional Gauss quadratures of a single disc--rectangle area, and the free-surface
condition ($\Delta p = 0$, no mass or viscous flux through the moving boundary) is imposed at the
true sphere $r=R$ by simply omitting the boundary face from the aperture-weighted finite-volume
divergence. Replacing the conventional stair-stepped mask, which imposes a rigid wall, by this
treatment moves the $\ell=2$ $f$-mode of the benchmark $\Gamma=2$ polytrope from $-6.95\%$ to
$-0.08\%$ of the frequency-domain shooting value at $N=48$ cells across the box, with $p_1$ to
$-1\%$; three smooth barotropes agree to $\le0.08\%$. With the cut-cell geometry made exact
(every quadrature split at its kinks and symmetrized over the axes), the engine seeds any real
$Y_{\ell m}$ on the full cube: the $E_g$ partners $Y_{20}$ and $\mathrm{Re}\,Y_{22}$ give
frequencies identical to fourteen significant figures, as the cubic group requires, while the
$T_{2g}$ member $\mathrm{Re}\,Y_{21}$ sits $1.13\%$ lower at $N=24$---a direct measurement of the
grid's cubic anisotropy. We extend the method to
Bemfica--Disconzi--Noronha--Kovtun (\BDNK) first-order causal viscous hydrodynamics, keeping the
frame relaxation times so that the linearized system remains hyperbolic, with the covariant
expansion and the lapse in both recovery equations. In the three hydrodynamic frames of
Clarisse~\emph{et al.} the viscous damping rate of the $f$-mode is frame-invariant to $0.045\%$
and converges to $99\%$ of the Navier--Stokes dissipation integral; the viscous frequency shift,
once the Kreiss--Oliger artifact that dominated it is removed, converges to the parameter-free
damped-oscillator value $-(\gamma^2-\gamma_0^2)/2\omega^2$ ($-0.125\%$ measured against
$-0.133\%$ predicted at $N=48$); while
the linear stability bound $\tau_Qw_0\cs^2 > \tfrac43\eta$ is shown to be sufficient but not
necessary: the marginally causal frame violates it in every cell yet evolves stably for 46
periods, because the unstable band sits above a critical wavenumber the stencil cannot resolve,
and a frame violating it in half the star destabilizes first. We report the method's limits with
equal weight: realistic crusted equations of state excite a growing surface mode for
$N\ge40$ whose discriminant is the surface index $n$ in $\varepsilon\propto(R-r)^n$; the
Kreiss--Oliger coefficient is a $\pm0.25\%$ systematic on absolute frequencies and the dominant
one on viscous frequency shifts; and time-domain frequency resolution at affordable cost is $4\%$
of $f$, so every sub-percent figure is a parametric estimate.
\end{abstract}

\maketitle

%======================================================================
\section{Introduction}
%======================================================================

Frequency-domain eigensolvers on spherically symmetric backgrounds compute neutron-star
oscillation modes to five or six figures and remain the standard of accuracy. A time-domain
3+1D evolution is worse at that task by construction: its spectral resolution is $1/T$, its
surface is not aligned with the grid, and its dissipation is not zero. The case for building one
anyway is what it can do that the eigensolvers cannot---carry a non-axisymmetric perturbation,
couple modes nonlinearly, and eventually sit inside a dynamical-spacetime merger code---and the
precondition for any of that is that the linear, axisymmetric, Cowling problem be reproduced
first. This paper is about that precondition.

The difficulty is the surface. A relativistic star's density and pressure go to zero at $r=R$,
so a Cartesian grid intersects a sphere of vanishing enthalpy along a stair-stepped set of cell
faces. The common treatments are an atmosphere (which introduces a spurious low-density fluid),
an excised shell (which imposes a wall inside the star), or a binary interior mask (which
imposes a wall on the staircase). We show below that the last of these---the most natural
choice for a linearized code, since exterior cells simply stay at zero---produces a $-7\%$
$f$-mode frequency error that does \emph{not} converge away between $N=32$ and $N=40$, from a
run that looks perfectly healthy. The error is not the geometric volume error of the staircase,
which is at the $10^{-4}$ level by $N=72$; it is the boundary condition. A stationary exterior
imposes $\dS_i = 0$ on the staircase, i.e.\ a rigid wall, where the physics of the $f$-mode
demands a free surface that moves with the fluid. Cut cells fix the condition, not the geometry.

The second subject is viscosity. The \BDNK\ formulation~\cite{BDN2022,Kovtun2019} is the first
first-order viscous theory that is causal, linearly stable and strongly hyperbolic, at the price
of a set of frame relaxation times $(\tau_\varepsilon,\tau_P,\tau_Q)$ on which no physical
prediction may depend. Clarisse~\emph{et al.}~\cite{Clarisse2025} have recently tested that
frame independence in a flux-conservative 1+1D code by running three distinct frames. We do the
same test on a star in 3+1D, and use the stellar setting to probe the linear stability bound
that the theory's own analysis provides.

Section~\ref{sec:model} states the equations; Sec.~\ref{sec:method} the discretization,
including the cut-cell geometry and the frequency-extraction protocol;
Sec.~\ref{sec:validation} validates the ideal engine against a shooting solver and the
literature; Sec.~\ref{sec:bdnk} reports the viscous results; Sec.~\ref{sec:dg} compares with a
nonlinear discontinuous-Galerkin engine on the same problem; Sec.~\ref{sec:limits} records the
method's limits and the defects found. Throughout, $G=c=1$; the benchmark star is the $n=1$,
$K=100$ polytrope with $\rho_c = 1.28\times10^{-3}$, $M = 1.4002\,\Msun$, $R = 14.154$~km, whose
relativistic-Cowling $\ell=2$ spectrum is tabulated in Ref.~\cite{Font2000} (Table~3) and
reproduced in Ref.~\cite{Nouri2022} (Table~VI) as $f, p_1, \ldots, p_4 = 1.8825, 4.1060,
6.0298, 7.8670, 9.6663\kHz$.

%======================================================================
\section{Physical model}
\label{sec:model}
%======================================================================

\subsection{Background in areal-Cartesian coordinates}
The TOV star $m' = 4\pi r^2\varepsilon$, $\nu' = 2(m+4\pi r^3p)/[r(r-2m)]$,
$p' = -(\varepsilon+p)(m+4\pi r^3p)/[r(r-2m)]$ is integrated by RK4 and embedded in a cube
$[-L,L]^3$, $L = 1.2R$, with $N$ cells per dimension. At each cell centre the 1D profile is
interpolated in $r = \sqrt{x^2+y^2+z^2}$. The areal metric
$ds^2 = -e^{\nu}dt^2 + e^{2\lambda}dr^2 + r^2d\Omega^2$, $e^{2\lambda} = (1-2m/r)^{-1}$,
becomes in Cartesian components
\begin{equation}
\gamma_{ij} = \delta_{ij} + (e^{2\lambda}-1)\,n_in_j, \qquad
\gamma^{ij} = \delta_{ij} + (e^{-2\lambda}-1)\,n_in_j,
\end{equation}
with $n_i = x_i/r$, $\sg = e^{\lambda}$ and lapse $\alpha = e^{\nu/2}$; the spatial metric has
eigenvalues $(e^{2\lambda},1,1)$, a purely radial stretch. Outside $r=R$ the Schwarzschild
exterior is used. The local gravity is $\Phi' = \nu'/2 = (m+4\pi r^3p)/[r(r-2m)]$ and
$p_0' = -w_0\Phi'$, $w_0 \equiv \varepsilon_0 + p_0$.

\subsection{Linearized Cowling equations in momentum form}
\label{sec:cowling}
Evolving the velocity perturbation $\dv^i$ fails at the surface: every term carries
$1/(\varepsilon_0+p_0)$. We evolve instead the covariant momentum density
$\dS_i = (\varepsilon_0+p_0)\dv_i$, in which that factor cancels identically from every
right-hand side (the Valencia momentum form applied to the linear problem). With the barotropic
closure $\dpp = \cs^2\de$ and $\dS^i = \gamma^{ij}\dS_j$,
\begin{align}
\partial_t\de &= -\frac{1}{\sg}\partial_i\big[\alpha\sg\,\dS^i\big] - \alpha\Phi'(n_i\dS^i),
\label{eq:cont}\\
\partial_t\dS_i &= -\alpha\big[\partial_i\dpp + \Phi'n_i(\dpp+\de)\big].
\label{eq:euler}
\end{align}
The lapse multiplying the whole of Eq.~\eqref{eq:euler} is the load-bearing detail. It follows
from $\nabla_\mu\delta T^{\mu}{}_\nu = 0$ written as
$\nabla_\mu T^{\mu}{}_\nu = (1/\sqrt{-g})\partial_\mu(\sqrt{-g}\,T^{\mu}{}_\nu) - \tfrac12T^{ab}\partial_\nu g_{ab}$:
the flux is $-(1/\sg)\partial_i(\alpha\sg\,\dpp)$, the source is
$\alpha[-\de\Phi' + \dpp\,\partial_i\ln\sg]$, and the $\partial_i\ln\sg$ pieces cancel between
them, leaving a \emph{bare} $\alpha$ on the flat-form right-hand side---not the
$\alpha\sg$-inside-the-divergence form, which leaves an uncancelled $\dpp\,\partial\ln\sg$ and
overshoots. The $\nu=t$ component of the same identity reproduces Eq.~\eqref{eq:cont} with zero
residual, which validates the framework. Without the lapse the continuity and Euler equations
carry different time normalizations and every mode frequency acquires a
resolution-\emph{independent} error of $\approx+5\%$. As an independent check, reducing
Eqs.~\eqref{eq:cont}--\eqref{eq:euler} to a single second-order radial ODE and changing dependent
variable, $E = \phi(r)V$ with $\phi'/\phi = -[(1+3\cs^2)\alpha' + \alpha\cs^{2\prime}]/(\alpha\cs^2)$,
carries them \emph{exactly} onto the McDermott--Van~Horn--Scholl $(W,V)$ system of the shooting
solver~\cite{Sotani2024}; without the lapse the map does not close. The two codes therefore
solve the same continuum problem and any residual is numerical.

The Cowling $f$-mode is undamped. The physical surface condition is the vanishing of the
Lagrangian pressure perturbation, $\Delta p = \dpp + \xi^rp_0' = 0$; for a polytrope
$p_0'(R) = 0$ so $\dpp(R) = 0$, and there is no mass flux and no viscous traction through the
moving surface.

\subsection{\BDNK\ causal viscous extension}
\label{sec:bdnkmodel}
The general-frame stress tensor~\cite{BDN2022} is
\begin{align}
T^{\mu\nu} &= (\varepsilon+\mathcal A)u^\mu u^\nu + (P+\Pi)\Delta^{\mu\nu}\notag\\
 &\quad - 2\eta\sigma^{\mu\nu} + u^\mu\mathcal Q^\nu + u^\nu\mathcal Q^\mu,\notag\\
\mathcal A &= \tau_\varepsilon[D\varepsilon + w\theta],\qquad
\Pi = -\zeta\theta + \tau_P[D\varepsilon + w\theta],\notag\\
\mathcal Q^\nu &= \tau_Q[wDu^\nu + \cs^2\Delta^{\nu\lambda}\partial_\lambda\varepsilon],
\label{eq:bdnk}
\end{align}
$D = u\cdot\nabla$, $\theta = \nabla\cdot u$, $w = \varepsilon+P$. Linearized about the static
star with $u^\mu_0 = (1/\alpha,0,0,0)$, the conserved densities carry \emph{proper}-time
derivatives of the primitives, $D\de = \alpha^{-1}\partial_t\de$, $(Du^i)_{\rm lin} =
\alpha^{-1}\partial_t\dv^i$, so the recovery is diagonal and each recovered derivative is
converted to coordinate time by a factor $\alpha$:
\begin{align}
\delta E &= \de + \tau_\varepsilon\big(\alpha^{-1}\partial_t\de + w_0\theta\big),\notag\\
\partial_t\de &= \alpha\big[(\delta E-\de)/\tau_\varepsilon - w_0\theta\big],
\label{eq:recE}\\
\dS^i &= w_0\dv^i + (\tau_Qw_0-\eta)\,\alpha^{-1}\partial_t\dv^i + \tau_Q\cs^2\gamma^{il}\partial_l\de,\notag\\
\partial_t\dv^i &= \frac{\alpha\big(\dS^i - w_0\dv^i - \tau_Q\cs^2\gamma^{il}\partial_l\de\big)}{\tau_Qw_0-\eta}.
\label{eq:recS}
\end{align}
Three details are exact on the curved background and each was found by failing without it. The
expansion is the four-dimensional one and carries both the lapse and $\sg$,
$\theta = (\alpha\sg)^{-1}\partial_i(\alpha\sg\,\dv^i) = \partial_i\dv^i + (\Phi'+A'/A)(n\cdot\dv)$,
not the flat divergence. The shear $\sigma^{ij}$ uses the closed-form Christoffels of
$\gamma_{ij} = \delta_{ij}+(A^2-1)n_in_j$, $\Gamma^a_{bc} = n^a[g_1n_bn_c + g_2(\delta_{bc}-n_bn_c)]$,
verified in all 27 components. And the bulk closure is pinned: $\Pi = \cs^2\mathcal A$ forces
$\tau_P = \cs^2\tau_\varepsilon$, and using $\tau_P = \tau_\varepsilon$ instead reopens a
primitive-feedback instability. The recovery denominator $\tau_Qw_0-\eta > 0$ is the
causality margin (BDN condition~(a), $\rho\tau_Q > \eta$). Because the frame times are kept,
the linearized system is hyperbolic and the diffusive CFL restriction of the Navier--Stokes
limit---unstable for $\hat\eta\gtrsim1$---does not arise.

A local analysis of the $4\times4$ longitudinal sound sector $(\delta E,\dS,\de,\dv)$ that this
system integrates, in flat space with frozen coefficients, gives the linear stability condition
\begin{equation}
\tau_Q\,w_0\,\cs^2 \;>\; \tfrac43\,\eta,
\label{eq:bound}
\end{equation}
\emph{not} $\tau_Qw_0 > \eta$; it was reproduced to ratio $1.0000$ in six $(\cs^2,\hat\eta)$
combinations and is used in Sec.~\ref{sec:bdnk}.

\subsection{Hydrodynamic frames}
Clarisse~\emph{et al.}~\cite{Clarisse2025} parametrize $\tau_\varepsilon = a_1(1/T)(\eta/s)$,
$\tau_Q = a_2(1/T)(\eta/s)$ under $a_1\ge4$, $a_2\ge3a_1/(a_1-1)$, with
$F_1 = (25/4,25/7)$, $F_2 = (25/2,25/3)$, $F_3 = (25,25)$ and characteristic speeds
$c_+ = 1.00, 0.85, 0.74$. $F_1$ saturates the second bound exactly, $3(25/4)/(21/4) = 25/7$,
which is why $c_+ = 1$: it is the marginally causal frame. At zero chemical potential
$w = sT$, so $(1/T)(\eta/s) = \eta/w$ and the stellar map is
\begin{equation}
\tau_\varepsilon(\bm x) = a_1\,\eta(\bm x)/w_0(\bm x),\qquad \tau_Q(\bm x) = a_2\,\eta(\bm x)/w_0(\bm x),
\label{eq:map}
\end{equation}
installed cell by cell, with $\eta = \hat\eta\,\varepsilon_0$. Under Eq.~\eqref{eq:map} the bound
\eqref{eq:bound} collapses to the $\eta$-independent $a_2\cs^2 > 4/3$. The repository
reproduces the published $c_+$ from $(a_1,a_2)$ alone via the Kovtun large-$k$
biquadratic~\cite{Kovtun2019}: $1.000000/0.849719/0.741915$, differing from the table by
two-decimal rounding.

%======================================================================
\section{Numerical method}
\label{sec:method}
%======================================================================

\subsection{Discretization}
Second-order central differences on the $N^3$ grid, classical RK4 with $\Delta t = 0.30\,\Delta x$
(ideal) or $0.20\,\Delta x$ (\BDNK), and fourth-derivative Kreiss--Oliger dissipation
$-\KO\sum_d(\ldots)$ applied to every evolved field. Exterior cells are never written; interior
cells are those with positive volume fraction (cut) or centre inside $R$ (masked). The engine is
single-threaded; cost scales as $N^4$ with a fitted $6.9\times10^{-8}$~s per $N^4\,\Msun$ of
evolution, i.e.\ $\sim900$~s for the production unit $N=48$, $T=2500\,\Msun$ (23 $f$-periods).

\subsection{Cut-cell geometry}
\label{sec:cutcell}
For each cell whose centre lies within $2\Delta x$ of $r=R$ we need the volume fraction
$\kappa$ inside the sphere and the six face apertures $a_{x\pm},a_{y\pm},a_{z\pm}$ (the fraction
of each face inside the sphere). Both reduce to one function: the area of the intersection of a
disc of radius $a$ with an axis-aligned rectangle $[y_0,y_1]\times[z_0,z_1]$,
\begin{align}
\mathcal D &= \int_{y_-}^{y_+}\!\ell(y)\,dy,\qquad y_\mp = \max(y_0,-a),\ \min(y_1,a),\notag\\
\ell(y) &= \big[\min(z_1,\zeta) - \max(z_0,-\zeta)\big]_+,\quad \zeta = \sqrt{a^2-y^2},
\end{align}
A face at $x = x_f$ meets the sphere in a disc of radius $\sqrt{R^2-x_f^2}$, so
$a_{x\pm} = \mathcal D(\sqrt{R^2-x_f^2};\ldots)/\Delta x^2$ directly; the volume fraction is a
second quadrature of $\mathcal D$ over $x$ across the cell, $\kappa = \Delta x^{-3}\int\mathcal D\,dx$.
No interface-plane reconstruction is needed.

The integrands are only piecewise smooth: $\ell(y)$ has kinks where the disc rim meets a
rectangle edge, $\mathcal D(a(x))$ has kinks where the rim passes an edge or a corner, and both
have a square-root singularity at the rim. A single Gauss rule therefore converges
algebraically---24 points left $\sim10^{-5}$ errors---and because $\kappa$ integrates over $x$
while the faces integrate over $y$ and $z$, that error was not symmetric under $x\leftrightarrow y$:
the geometry broke the sphere's own octahedral symmetry at $10^{-5}$, which leaked into the
$Y_{\ell m}$ moments (Sec.~\ref{sec:fullcube}). The implementation is therefore exact: the
substitutions $y = a\sin\theta$ and $x = R\sin\varphi$ remove the rim singularities, every
integral is split at its kink points ($\cos\theta = \{z_1,-z_0,z_0,-z_1\}/a$;
$\cos\varphi = \{|y_{0,1}|,|z_{0,1}|,\text{corner distances}\}/R$), and 16-point Gauss on each
smooth piece is exact to round-off; $\kappa$ is finally averaged over the six axis permutations,
under which the true $\kappa$ is invariant. Verified against the analytic sphere:
$\sum\kappa\,\Delta x^3 = \tfrac43\pi R^3$ to $2\times10^{-16}$, the $x$-face apertures equal the
transposed $y$-face apertures to $0.0$, and the closure condition---the face-aperture area
vectors balancing the boundary face---reproduces $4\pi R^2$ to $8\times10^{-4}$ at $N=24$ and
$4.5\times10^{-4}$ at $N=32$ (second order). Cells with $\kappa$ below a floor
$\kappa_{\min} = 0.5$ use $\kappa_{\min}$ in the denominator (a small-cell treatment).

\subsection{Aperture-weighted divergences and the free surface}
Every spatial derivative in Eqs.~\eqref{eq:cont}--\eqref{eq:euler} is a divergence or a gradient,
and both are evaluated in finite-volume form. For the continuity flux $F^i = \alpha\sg\,\dS^i$,
\begin{align}
(\nabla\!\cdot\!F)_{ijk} = \frac{1}{\kappa_c\Delta x}\Big[&a_{x+}\tfrac12(F^x_{i}+F^x_{i+1})\notag\\
 &- a_{x-}\tfrac12(F^x_{i-1}+F^x_i) + (y) + (z)\Big],
\label{eq:divF}
\end{align}
with face values the second-order average of the adjacent cell-centred fluxes. The gradient of
$\dpp$ uses the divergence theorem in the same form, $\int\partial_j\dpp\,dV = \oint\dpp\,n_j\,dA$,
so coordinate faces contribute $a\,\dpp$; and the viscous force $\partial_j\pi^{ij}$ likewise.
\emph{The boundary face does not appear in any of these sums.} That omission \emph{is} the
free-surface boundary condition: the boundary face would contribute $\dS^i n_i A_b$ (mass flux),
$\dpp_bn_jA_b$ (pressure) and $\pi^{ij}n_jA_b$ (traction), all of which vanish for a free
surface that moves with the fluid and carries $\Delta p = 0$. The condition is thereby imposed at
the true sphere, and the exterior is never referenced. In the masked engine, by contrast, the
central stencil just inside the staircase reads the frozen exterior zeros, imposing $\de = 0$
(correct, since $\dpp(R)=0$) \emph{and} $\dS_i = 0$---a rigid wall. This is the whole difference,
and Sec.~\ref{sec:validation} quantifies it.

Two remedies that do \emph{not} work are recorded so they are not retried. Zero-gradient
extrapolation of $\dS$ into an exterior collar with $\de = 0$ makes the mode collapse
($-61\%$ at $N=32$, $-83\%$ at $N=44$): momentum in a region with no pressure restoring force
plus a spurious mass flux in the continuity stencil. And an excised shell (evolving only
$w_0 \ge f_{\rm exc}\max w_0$) merely moves the wall inward, giving $+2.6\%$ at $f_{\rm exc}=0.02$
and $+8.2\%$ at $0.08$.

\subsection{Seeding and diagnostics}
The seed is a density perturbation $\de = A\,\varepsilon_0\,(r/R)^\ell\,Y_{\ell m}(\hat n)$ with
$Y_{\ell m}$ the real solid-harmonic angular factor ($\ell = 2$--$4$, all $m$; cosine type for
$m\ge0$, sine type for $m<0$) and $A = 10^{-3}$; the radial factor makes $r^\ell Y_{\ell m}$ a
polynomial, smooth at the centre. The legacy $\ell=2$, $m=0$ seed used $(r/R)^1$ and is retained.
The diagnostics are the $\sg$-weighted moments $q_{\ell m}(t) = \sum\de\,Y_{\ell m}(\hat n)\,\sg\,
\Delta x^3$ over interior cells, any set of which can be recorded in one run. Because the code is
linear, amplitude proportionality is exact to $10^{-14}$ (verified over $\times100$ and
$\times3000$), which makes the linear engine the control for any nonlinear-engine claim
(Sec.~\ref{sec:dg}). The grid is the full cube, not an octant, so every $m$ is admissible.

\subsection{Frequency extraction and its resolution}
\label{sec:extract}
Two independent estimators are applied to every record and must agree: a direct periodogram
$|\sum_nq_n e^{-i\omega t_n}|^2$ with parabolic refinement of the peak, and a matrix pencil
(Hankel SVD, 14--16 poles) on the record decimated to $\sim600$ samples. Every record is gated
before either is applied: finiteness, and the ratio of the envelope at the end to that at the
start; a growing record is re-analysed on its largest clean window, and if that window is
shorter than ten periods the row is voided rather than reported. A record of length $T$ has
spectral resolution $\Delta f = 1/T$, which at the production $T = 2500\,\Msun$ is
$0.083\kHz = 4.4\%$ of $f$. \emph{Every sub-percent figure below is therefore a parametric peak
estimate, quoted only where the two estimators agree, and no third decimal is meaningful.}
Buying record length does not help: at $\KO = 0.004$, $T = 12000$ actively degrades the answer as
a late-time surface mode takes over ($-0.30\%\to-1.67\%$). The damping rate is the slope of a
log-linear fit to the local-maximum envelope, returning \texttt{NaN} (not zero) when fewer than
three peaks exist.

The viscous damping is measured by a \emph{difference protocol}: every frame is run twice at
identical grid and frame times, once at $\hat\eta$ and once at $\hat\eta=0$, and
$\gamma_{\rm visc} = \gamma(\hat\eta)-\gamma(0)$. The numerical floor $\gamma(0)$ from
Kreiss--Oliger and the surface is $1.30\times10^{-3}\,\Msun^{-1}$ at $N=40$---$52\%$ of the
viscous rate---so quoting $\gamma(\hat\eta)$ raw would overstate the damping by $1.5\times$.

%======================================================================
\section{Validation of the ideal engine}
\label{sec:validation}
%======================================================================

\subsection{The benchmark star}
Table~\ref{tab:anchor} is the central validation. With cut cells at $N=48$ the $f$-mode is
$-0.083\%$ from the shooting value and $p_1$ is $-0.98\%$; the masked engine at the same
$\KO$ is $-6.95\%$ at both $N=32$ and $N=40$, i.e.\ the error does not converge in this range
and would be reported as a healthy result by any pipeline that did not compare to a reference.
The cut-cell sequence $N = 24/32/40/48$ is non-monotone ($+0.005/-0.39/-0.31/-0.08\%$), so no
convergence order can be fitted, Richardson extrapolation on (40,48) makes the answer
\emph{worse} ($+0.44\%$), and the honest error bar is the $N$-spread of $\pm0.4\%$ rather than
the $N=48$ residual. The Kreiss--Oliger coefficient is a systematic of the same size:
$\KO = 0.02$ gives $+0.42\%$ at $N=48$, bracketing the reference from above, and any cut-cell
frequency must state $\KO$; on the masked path it is worth more than seven points. $p_2$
($5.771$ vs $6.031\kHz$) is not resolved at these record lengths.

\begin{table}[tb]
\footnotesize
\caption{Benchmark $\Gamma=2$ polytrope, $\ell=2$, $m=0$, $T=2500\,\Msun$ (23 periods,
$\Delta f/f = 4.4\%$), $\kappa_{\min}=0.5$. Reference: shooting solver, $f = 1.88291$,
$p_1 = 4.10672\kHz$. Periodogram values; matrix pencil agrees to $\le0.13\%$ in every row.}
\label{tab:anchor}
\begin{ruledtabular}
\begin{tabular}{lcclcc}
surface & $N$ & $\KO$ & mode & $f$ [kHz] & dev\\
\hline
cut    & 24 & 0.01 & $f$ & 1.88301 & $+0.005\%$\\
cut    & 32 & 0.01 & $f$ & 1.87560 & $-0.388\%$\\
cut    & 40 & 0.01 & $f$ & 1.87701 & $-0.313\%$\\
cut    & 48 & 0.01 & $f$ & \textbf{1.88135} & $\bm{-0.083\%}$\\
cut    & 48 & 0.01 & $p_1$ & 4.0665 & $-0.98\%$\\
cut    & 48 & 0.02 & $f$ & 1.89076 & $+0.417\%$\\
masked & 32 & 0.004 & $f$ & 1.75235 & $-6.93\%$\\
masked & 40 & 0.004 & $f$ & 1.75205 & $-6.95\%$\\
\end{tabular}
\end{ruledtabular}
\end{table}

The literature target is the \emph{linear} Cowling column of Font~\emph{et al.}~\cite{Font2000}
(Table~3, Model~2), $1.8825\kHz$, which our shooting solver reproduces to $0.02\%$. Font's
same table lists a nonlinear HRSC evolution on a fixed spacetime at $1.852\kHz$; the $1.7\%$
gap is a finite-amplitude effect our linearized code must not reproduce, and the full-GR value
for this star, $1.578$--$1.586\kHz$~\cite{KrugerKokkotas2020,Nouri2022}, is $19\%$ below
Cowling and is not a valid comparison for a frozen-metric engine.

\subsection{Non-axisymmetric seeds on the full cube}
\label{sec:fullcube}
Under the three reflections $x\to-x$, $y\to-y$, $z\to-z$ of the cubic group the $\ell=2$ quintet
splits into the irreps $E_g = \{Y_{20},\mathrm{Re}\,Y_{22}\}$ (even under all three) and
$T_{2g} = \{\mathrm{Re}\,Y_{21},\mathrm{Im}\,Y_{21},\mathrm{Im}\,Y_{22}\}$ (odd under two). Two
statements follow for any discretization that respects the cube's symmetry: cross-moments between
irreps vanish identically, and the two $E_g$ partners cannot be split at all. Both are exact
tests of the discretization rather than of the physics, and both are met. Seeding $Y_{20}$,
$\mathrm{Re}\,Y_{22}$ and $\mathrm{Re}\,Y_{21}$ in turn at $N=24$ and recording all three moments
in each run, every cross-moment is below $10^{-8}$ of the seeded one, and the $E_g$ partners give
$f = 1.884\,603\,501\,533\,21$ and $1.884\,603\,501\,533\,23\kHz$---identical to fourteen
figures. With the previous 24-point Gauss geometry the same pair differed at $5\times10^{-7}$ and
the cross-moments were $10^{-5}$: the asymmetric $\kappa$ was the entire effect. The $T_{2g}$
member, by contrast, is genuinely split from $E_g$ by the grid: $f(\mathrm{Re}\,Y_{21}) =
1.8633\kHz$, $-1.13\%$, at $\Delta x/R = 0.10$. This is the cubic anisotropy of the second-order
stencil made visible as a mode splitting; on a spherical star it must vanish as $h^2$, and it is
a cleaner measurement of it than the nonlinear engine could give (Sec.~\ref{sec:dg}).

\subsection{Across equations of state}
\label{sec:eos}
The background builder accepts any barotrope, and the Read~\emph{et al.} piecewise polytropes
SLy, APR4, H4 and MS1~\cite{Read2009} reproduce their 1D masses and radii to six figures on the
grid. Table~\ref{tab:eos} shows the outcome. The three smooth analytic barotropes converge to
$\le0.08\%$ at $N=48$, all dipping to $-0.3$ to $-0.7\%$ at intermediate $N$ first. The four
crusted EOS behave differently: at $\KO = 0.01$ a surface mode grows at $+0.7$ to $+1.6$
e-folds per $f$-period for $N\ge40$ (end/start amplitude ratios $4\times10^{11}$ to
$8\times10^{14}$ over 23 periods), so they are usable only at $N\le32$, at $-1.6$ to $-4\%$;
Richardson on the monotone (24,32) pair gives $-0.5\%$ (APR4) to $-2.6\%$ (H4).

\begin{table*}[tb]
\caption{Cut-cell $f$-mode across EOS at $M\simeq1.4\,\Msun$, $\KO=0.01$ unless stated,
against the shooting value for the identical star. R = Richardson on the monotone (24,32) pair.}
\label{tab:eos}
\begin{ruledtabular}
\begin{tabular}{llccl}
EOS ($M/\Msun$, $R$/km) & $N$ & $f_{3D}$ & $f_{1D}$ & dev\\
\hline
$\Gamma=2$ (1.400, 14.15) & 48 & 1.88135 & 1.88291 & $-0.083\%$\\
$\Gamma=2.5$ & 48 & 2.23448 & 2.23404 & $+0.020\%$\\
$\Gamma=3$   & 48 & 2.40545 & 2.40739 & $-0.080\%$\\
SLy (1.421, 11.67) & 32 & 2.37154 & 2.41928 & $-1.97\%$; R $-1.02\%$\\
 & 40, 48 & \void & & surface mode\\
 & 48, $\KO{=}0.04$ & 2.31598 & & $-4.27\%$ (biased)\\
APR4 (1.360, 11.31) & 32 & 2.45373 & 2.49463 & $-1.64\%$; R $-0.51\%$\\
 & 40, 48 & \void & & end/start $8\times10^{14}$\\
H4 (1.406, 13.93) & 32 & 1.89169 & 1.94832 & $-2.91\%$; R $-2.57\%$\\
MS1 (1.359, 14.90) & 24 & 1.69672 & 1.76756 & $-4.01\%$ (only stable $N$)\\
 & 40, 48, $\KO{=}0.04$ & 1.646, 1.643 & & $-6.9$, $-7.1\%$\\
\end{tabular}
\end{ruledtabular}
\end{table*}

What discriminates the two classes is neither compactness, stiffness nor tabulation but the
steepness of the surface. Writing $\varepsilon\propto(R-r)^n$ between $0.99R$ and $0.999R$, the
effective index is $n = 0.50, 0.67, 1.005$ for $\Gamma = 3, 2.5, 2$ (stable) and
$1.84, 1.87, 1.95, 1.95$ for MS1, H4, APR4, SLy (unstable), and the ordering in $n$ is the ordering
in the $N=48$ end/start ratio, monotone over nineteen decades with no exception. The SLy crust's
outermost segment has $\Gamma\simeq1.58$, hence $n = 1/(\Gamma-1)\simeq1.7$--$2.0$, and
$\varepsilon_0+p_0$ falls by a factor $\sim500$ across one grid cell exactly where the apertures
are partial, against $0.44$--$0.70$ for the analytic EOS. This is a correlation across seven EOS,
not a demonstrated mechanism; the one causal hypothesis tested---that the finite atmosphere
sound speed of the tabulated EOS ($\cs^2\sim3\times10^{-5}$ versus $10^{-10}$--$10^{-19}$) drives
it---was refuted directly by zeroing $\cs^2$ in all $77\,040$ atmosphere cells of the SLy
$N=48$ run: growth $+0.596$ versus $+0.595$ e-folds per period. Raising $\KO$ to $0.04$
stabilizes every crusted run for the full 23 periods but biases $f$ by $-3.8$ to $-7.1\%$, with the
bias \emph{growing} with $N$; there is no configuration we could afford that gives a converging
crusted-EOS $f$-mode, and the mechanism is open.

%======================================================================
\section{Causal viscous results}
\label{sec:bdnk}
%======================================================================

\subsection{Frame robustness of the damped $f$-mode}
Table~\ref{tab:frames} gives the $\ell=2$ $f$-mode at fixed physical viscosity $\hat\eta=0.03$,
$\zeta=0$, in the three Clarisse frames, with the $\hat\eta=0$ control and the difference
protocol. The frame times under Eq.~\eqref{eq:map} are nearly uniform across the star because
$\eta\propto\varepsilon_0$: $\tau_\varepsilon\in[0.168,0.188]$, $[0.337,0.375]$, $[0.674,0.750]\,\Msun$
for $F_{1,2,3}$.

\begin{table*}[tb]
\caption{\BDNK\ $f$-mode in the three frames, $\hat\eta=0.03$, cut cells, $\KO=0.01$, $N=40$,
$T=2500\,\Msun$ (22.8 periods). $\gamma_{\rm visc}=\gamma(\hat\eta)-\gamma(0)$; 1D polar values
from the $\ell$-reduced frequency-domain operator with the corrected free-surface BC.}
\label{tab:frames}
\begin{ruledtabular}
\begin{tabular}{lcccc}
 & $F_1$ & $F_2$ & $F_3$ & spread\\
\hline
$(a_1,a_2)$ & $(25/4,25/7)$ & $(25/2,25/3)$ & $(25,25)$ & \\
$f$ pencil [kHz] & 1.87954 & 1.87973 & 1.88040 & $0.046\%$\\
$f$ periodogram & 1.88027 & 1.88036 & 1.88082 & $0.029\%$\\
$f$, 1D polar & 1.9038 & 1.9105 & 1.9272 & $1.22\%$\\
$f$, $\hat\eta=0$ control & 1.87681 & 1.87681 & 1.87681 & $0$\\
$\gamma_{\rm visc}$ [$10^{-3}\Msun^{-1}$] & 2.48251 & 2.48229 & 2.48139 & $\bm{0.045\%}$\\
$\tau_{\rm damp}$, $Q$ & 1.984 ms, 11.72 & 1.984 ms, 11.72 & 1.985 ms, 11.73 & \\
cells violating \eqref{eq:bound} & $100\%$ & $94.4\%$ & $51.3\%$ & \\
stable, 22.8 periods & yes & yes & yes & \\
\end{tabular}
\end{ruledtabular}
\end{table*}

The frame spread of the damping rate shrinks with resolution, $0.150\to0.130\to0.045\%$ at
$N=24/32/40$, the signature of a frame-invariant physical rate whose apparent spread is
discretization residue. The rate itself converges at observed order $\approx3$
($2.35482\to2.44941\to2.48251\times10^{-3}$; successive differences in ratio $2.86$), and
Richardson extrapolation gives $2.51$--$2.54\times10^{-3}\,\Msun^{-1}$. That sequence is at
$\KO = 0.01$; at $\KO = 0.005$ the matrix-pencil rate is $2.597\times10^{-3}$ at $N=40$ and
$2.607\times10^{-3}$ at $N=48$---$99.2\%$ of the Newtonian-form Navier--Stokes dissipation
integral $0.08746\,\hat\eta = 2.6238\times10^{-3}$ for this star. The apparent few-percent
``causal suppression'' of the $\KO=0.01$ sequence was Kreiss--Oliger bias. The frame spread of
the frequency, $<0.05\%$, is an upper bound only ($\Delta f/f = 4.4\%$), but is $25\times$ smaller
than the $1.22\%$ of the 1D polar operator; that operator's own damping rates were withheld
after its velocity recovery was found to omit the lapse, and the 3D engine---which carries the
lapse in Eq.~\eqref{eq:recS}---independently confirms the retraction. The viscous frequency
shift relative to the inviscid control is the one quantity that required a dedicated
convergence study (Sec.~\ref{sec:shift}).

\subsection{$a_1$ is inert}
Holding $a_2 = 25$ and sweeping $a_1 = 6.25/25/100$ gives $f = 1.89112/1.89113/1.89123\kHz$
(spread $0.006\%$) and $\gamma = 5.618/5.621/5.630\times10^{-3}$ (spread $0.22\%$) over a
$16\times$ range. The three Clarisse frames differ in the stellar polar sector through $a_2$
alone. In the axial sector this is exact---$\mathcal A$ is a scalar and odd-parity perturbations
carry no scalar content---and here it is empirical.

\subsection{The viscous frequency shift}
\label{sec:shift}
At $\KO = 0.01$ the shift $\Delta f/f = f(\hat\eta)/f(0) - 1$ measured $+0.40/+0.28/+0.15\%$ at
$N = 24/32/40$ with non-shrinking differences, and could only be quoted as an upper bound. Two
observations settled it (33 runs; \texttt{repro/bdnk3d\_viscous\_shift.jl}). First, with the
frame times held fixed and $\hat\eta$ swept, the shift is roughly linear in $\hat\eta$, whereas
every physical mechanism---the damped-oscillator correction and the reactive frame terms---is
second order. Second, halving $\KO$ removes $\approx0.35\%$ of the shift at every $N$: a physical
shift cannot depend on the dissipation coefficient, so the positive shift was a Kreiss--Oliger
artifact. What survives at $\KO = 0.005$ converges to the damped-oscillator value, which has no
free parameter because both damping rates are measured in the same runs:
\begin{equation}
\frac{\Delta f}{f} = -\frac{\gamma_{\rm tot}^2 - \gamma_0^2}{2\omega^2}.
\label{eq:dosc}
\end{equation}
At $\hat\eta = 0.03$ the measured shift (mean of the two estimators) against
Eq.~\eqref{eq:dosc} is $+0.05/-0.20\%$, $-0.03/-0.16\%$, $-0.08/-0.13\%$ and
$-0.125/-0.133\%$ at $N = 24/32/40/48$; the residual $+0.25\to+0.13\to+0.05\to+0.01\%$ falls
roughly as $h^3$, and at $N=40$ the $\hat\eta$-dependence is quadratic once that offset is
removed ($\hat\eta = 0.015, 0.03, 0.045$). Frame-reactive corrections are $O(\tau^2\omega^2)
\approx 3\times10^{-5}$, below resolution. The converged statement is
$\Delta f/f(\hat\eta{=}0.03) = -0.125\pm0.013\%$, the damped-oscillator shift. Three protocol
facts the data enforce: the frame must be held fixed while $\hat\eta$ is swept, or the
$\eta$-dependence is entangled with $\tau\propto\hat\eta$; the inviscid control becomes unstable
below $\KO\approx0.005$ at $N = 24$ and $32$ (the cut-cell surface mode) while viscous runs
survive to $0.00125$, so $\KO = 0.005$ is the floor of the difference protocol and is set by the
control; and the compact viscous stencil (\texttt{visc\_compact}) shifts $f$ by $+0.15\%$ and
cuts $\gamma_{\rm visc}$ by $32\%$ at $N=24$ and is not to be used for the mode.

\subsection{The stability bound is sufficient, not necessary}
\label{sec:stability}
$F_1$ violates Eq.~\eqref{eq:bound} in $100\%$ of cells and $100\%$ of proper volume, yet
evolves stably for 46 periods through twelve decades of decay with a window-by-window decay rate
constant to $1\%$; it is as clean as $F_3$, which violates it in $51\%$. The local sound-sector
analysis explains why. The unstable band lies above a critical wavenumber
$k_c = C/\hat\eta$ with $C = 0.4599, 0.3019, 0.1753$ for $F_{1,2,3}$ (minimum over the star, always
at $r/R = 0.994$), i.e.\ $k_c = 15.3, 10.1, 5.8\,\Msun^{-1}$ at $\hat\eta = 0.03$. The
second-order centred stencil caps the effective wavenumber at $k_{\rm eff} = \sin(kh)/h \le 1/h =
N/23.0$, so the unstable band is resolved only for $N > 353, 232, 135$---beyond reach at
$\sim N^4$ cost. The criterion ``unstable iff $k_c < 1/h$'' was therefore tested on the cheap
$\hat\eta$ axis at $N=24$ ($1/h = 1.043$): six predictions, six correct in sign, with growth
rates agreeing to $24$--$28\%$ where growth is well developed ($F_1$, $\hat\eta=1.0$: predicted
$3.50\times10^{-2}$, measured $4.49\times10^{-2}$). And the ordering inverts: at $\hat\eta=0.60$,
$F_3$ blows up ($q_{20}$ grows $5\times10^4$) while $F_1$ decays, because $k_c(F_3)$ is
$2.6\times$ smaller. \emph{The violation fraction does not order the instability; the critical
wavenumber does.} The published fractions $60.7\%$ and $21.4\%$ from the 1D radial grid for
$F_2,F_3$ are a counting artifact; the 3D cell fractions match the 1D proper-\emph{volume}
fractions ($94.6\%$, $52.6\%$) to $<1\%$.

%======================================================================
\section{Comparison with a nonlinear DG engine}
\label{sec:dg}
%======================================================================

The same problem was run in a full 3+1D nodal Runge--Kutta discontinuous-Galerkin Valencia
GRHydro code on the frozen metric (densitized conserved variables, directional Rusanov sweeps,
SSP-RK3, TVB/Zhang--Shu limiter, atmosphere, well-balanced background subtraction), on an octant
grid with a full 3-momentum. It is the only engine here that can carry a non-axisymmetric
perturbation, and it does: a $\mathrm{Re}\,Y_{22}\propto n_x^2-n_y^2$ seed is detected by the
$m=2$ quadrupole with contrast $3\times10^5$--$5\times10^6$ over the $m=0$ seed at every
resolution. On a non-rotating star the $m$-degeneracy is exact, so the measured
$|f(m{=}2)-f(m{=}0)|/f$ of $6.4/3.8/2.1/0.81\%$ at $K = 5/6/8/10$ elements per axis is a grid
artifact falling as $h^{3.2}$, consistent with the scheme's $O(h^3)$; a single run seeded with
$Y_{21}$, whose mirror extension populates both $E_g$ channels on one background, bounds the
grid's genuine degeneracy violation at $9\times10^{-5}$---a bound that the linear engine's
fourteen-figure $E_g$ agreement (Sec.~\ref{sec:fullcube}) has since superseded by ten orders of
magnitude. Two facts, however, make the linear cut-cell method the one to use for precision. The DG background is not static: the well-balanced
subtraction is exact but the limiter is not, and $\rho_c$ erodes $-10\%$ by $4.6$ periods and
$-31$ to $-43\%$ by $3000\,\Msun$, capping the usable record at $\Delta f/f = 21.6\%$ and putting
its absolute $f$ at $-16\%$. And the octant's three separate reflections admit only the $E_g$
doublet $\{Y_{20},\mathrm{Re}\,Y_{22}\}$; the entire $T_{2g}$ triplet, including all of $m=1$, is
structurally absent from the domain. A search for nonlinear amplitude dependence in the DG
engine returned no resolved shift: the one suggestive point ($-1.4\pm1.0\%$ at seed amplitude
$0.1$, $K=8$) becomes $-7.5\pm0.9\%$ at $K=6$, a $5.4\times$ change for a $25\%$ refinement, and
no second harmonic is detected ($A_{2f}/A_f < 0.09$--$0.19$ at $\xi/R = 0.41$). The linear engine's
exact amplitude proportionality was the control that established this.

%======================================================================
\section{Limitations and defects}
\label{sec:limits}
%======================================================================

\emph{Structural.} The metric is frozen: no gravitational-wave damping, no $w$-modes, no tidal
deformability. The engine is linear. Its public interface is now a single call,
\texttt{bdnk3d\_qnm}, which builds the star and engine, installs a Clarisse frame from
$(a_1,a_2)$ or a scalar $\hat\tau$, seeds any $(\ell,m)$, evolves, applies the gated
two-estimator read-out and the $\hat\eta=0$ difference protocol, and returns the frequency,
damping, floor, Rayleigh resolution, stability flag and bound-violation fraction together;
\texttt{test/test\_bdnk3d\_qnm.jl} verifies each of these against the numbers in this paper at
$N=24$.

\emph{Numerical.} Time-domain resolution at affordable cost is $\Delta f/f = 4.4\%$; $\KO$ is a
$\pm0.25\%$ systematic on absolute frequencies and, at $\KO=0.01$, produced a spurious
$+0.3\%$ viscous frequency shift (Sec.~\ref{sec:shift}); the cut-cell path develops a late-time growing surface mode at
$\KO = 0.004$ for $N\ge32$ (growth $3\times10^{-4}$ to $9\times10^{-3}\,\Msun^{-1}$ from $N=32$
to $56$), cured at $\KO=0.01$ through $N=48$ but unverified above; realistic crusted EOS are
limited to $N\le32$; the numerical damping floor $\gamma(0)\approx1.2\times10^{-3}\,\Msun^{-1}$ at
$N=48$ is a hard floor under any viscous measurement without the difference protocol; and $p_2$
is not resolved.

\emph{Defects found in the code, all recorded rather than silently fixed.} The
\texttt{setup\_evo3d} default is \texttt{cut=false}, i.e.\ the $-6.95\%$ masked engine, with no
symptom. The background builder sets $\rho_0 = \varepsilon-p$, exact only for $\Gamma=2$; harmless
to these linear engines, which never read $\rho_0$, but fatal to any nonlinear user. The
$\ell$-reduced 1D polar operator's velocity recovery omits the lapse. And a Julia lexical trap:
\texttt{2f1} parses as the \texttt{Float32} literal $20.0$, not $2\,f_1$, which turned a
harmonic search into a fake null until a synthetic-injection test caught it.

%======================================================================
\section{Conclusions}
%======================================================================

A linearized Cowling evolution on a non-conforming Cartesian grid reproduces the frequency-domain
$f$-mode of smooth relativistic polytropes to better than $0.1\%$ once two things are right that
are easy to get wrong: the lapse on the Euler equation, and the boundary condition at the stellar
surface. The first is a derivation. The second is a cut-cell finite-volume divergence from which
the boundary face is simply absent---which imposes a free surface at the true sphere instead of a
rigid wall on the staircase, and is worth seven percent. Extended to \BDNK\ with the frame times
kept and the recovery written in proper time, the engine gives a viscous damping rate that is
frame-invariant to $0.045\%$, converges at third order to within a few percent of the
Navier--Stokes dissipation integral, and exposes the linear stability bound of the theory as
sufficient but not necessary on a star, with the critical wavenumber rather than the violation
fraction controlling onset. The crusted-EOS surface instability is the open problem; its
discriminant is known and its mechanism is not.

\begin{acknowledgments}
The validation runs, the frame study and the DG comparison were executed as adversarially
verified agent workflows: every headline number was re-run by an independent verifier, and the
verifiers' retractions are incorporated above.
\end{acknowledgments}

\begin{thebibliography}{99}
\bibitem{BDN2022} F.~S.~Bemfica, M.~M.~Disconzi and J.~Noronha, Phys.\ Rev.\ X \textbf{12},
021044 (2022).
\bibitem{Kovtun2019} P.~Kovtun, JHEP \textbf{10}, 034 (2019), arXiv:1907.08191.
\bibitem{Clarisse2025} N.~Clarisse, E.~O.~Pinho, T.~Patel, F.~S.~Bemfica, M.~Hippert and
J.~Noronha, ``Flux-conservative BDNK hydrodynamics and shock regularization,'' arXiv:2510.16603.
\bibitem{Font2000} J.~A.~Font, N.~Stergioulas and K.~D.~Kokkotas, Mon.\ Not.\ R.\ Astron.\ Soc.\
\textbf{313}, 678 (2000).
\bibitem{Nouri2022} F.~H.~Nouri, S.~Bose, M.~D.~Duez and A.~Das, arXiv:2107.13339.
\bibitem{KrugerKokkotas2020} C.~J.~Kr\"uger and K.~D.~Kokkotas, Phys.\ Rev.\ D \textbf{102},
064026 (2020), arXiv:2008.04127.
\bibitem{Sotani2024} H.~Sotani \emph{et al.}, arXiv:2412.18569.
\bibitem{Read2009} J.~S.~Read, B.~D.~Lackey, B.~J.~Owen and J.~L.~Friedman, Phys.\ Rev.\ D
\textbf{79}, 124032 (2009).
\end{thebibliography}

\end{document}
